\appendices
\section{Detailed Quantization and Constraint Reasoning}\label{app:bounds}
For the base objective, express each rounded coefficient as its exact scaled value plus an error in $[-1/2,1/2]$. For a fixed $k$-element set, the total absolute error is bounded by
\begin{equation}
\frac{1}{2}\left(k+\binom{k}{2}\right)=\frac{k(k+1)}{4}.
\end{equation}
Dividing by $A_bk$ yields~\eqref{eq:base_quantization}. Errors may cancel, so the actual error can be much smaller than this worst-case bound. The result concerns a fixed set, not the outputs of two different stochastic optimization trajectories.

For refinement, let $e_A(S)=E_r(S)/(28A)-\red(S)$. Equation~\eqref{eq:refinement_bound} states $|e_A(S)|\leq1/(2A)$. Exact optimality in the encoded objective implies
\begin{align}
\red(S_A)+e_A(S_A)&\leq\red(S_\star)+e_A(S_\star),\\
\red(S_A)-\red(S_\star)&\leq e_A(S_\star)-e_A(S_A)\leq1/A.
\end{align}
The admissible domain must be the same for both optima. The comparison does not cover a quantized optimizer using another pool, fixed head, or score constraint. A GAS output that misses the exact encoded optimum has an additional optimization gap.

The equivalence between scores and counts includes ties. For the ascending triple $y_1\leq y_2\leq y_3$, $y_j\geq\theta_j$ holds exactly when at least $4-j$ values are at least $\theta_j$. If $y_j<\theta_j$, the first $j$ entries fail the threshold, leaving at most $3-j$ that can meet it. The same ordering argument gives the reverse implication. Applying it to each threshold proves~\eqref{eq:threshold_counts} without requiring distinct logits.

\section{Native Validation Bound and Reporting Conventions}\label{app:validation}
After applying the independent expected-state inverse, the ideal joint decision/workspace state is all zero. If $p_j=(1-\langle Z_j\rangle)/2$ denotes the probability that qubit $j$ is one, a union bound gives
\begin{equation}
1-|\langle 0\cdots0|\psi\rangle|^2\leq\sum_j p_j.
\label{eq:validation_bound}
\end{equation}
Restricting the sum to workspace qubits bounds the probability of dirty workspace as well. Native validation checks the bound and norm against $10^{-8}$ for the small instances actually simulated. The check is phase-sensitive for those circuits; it does not validate the 42 refused simulations or other unrun instances.

Table~\ref{tab:thresholds} retains the saved legacy threshold analysis. Its no-marked percentage is calculated over queries with recorded cosine candidates; seeds are not treated as independent events. Development outcomes come from saved matched-sampling rankings, not a new threshold sweep on test judgments. A lower threshold can reduce empty marked sets but also change the ranking policy, making fallback rate insufficient as an effectiveness criterion.

\begin{table}[t]
\centering\small\setlength{\tabcolsep}{3pt}
\caption{Recorded legacy development thresholds. nDCG@5 comes from saved matched-sampling outcomes. No-marked percentages count development queries, with no new ranking or test tuning.}
\label{tab:thresholds}\begin{tabular}{lrrr}
\toprule
Dataset & Threshold & Dev. nDCG@5 & No-marked queries (\%)\\
\midrule
SciFact & 0.3 & 0.6627 & 0.0\\
SciFact & 0.5 & 0.6579 & 6.8\\
SciFact & 0.7 & 0.6656 & 67.1\\
NFCorpus & 0.3 & 0.3010 & 0.6\\
NFCorpus & 0.5 & 0.3011 & 31.8\\
NFCorpus & 0.7 & 0.3026 & 87.0\\
FiQA & 0.3 & 0.3514 & 0.0\\
FiQA & 0.5 & 0.3515 & 1.8\\
FiQA & 0.7 & 0.3516 & 69.4\\
\bottomrule
\end{tabular}

\end{table}

The complete exports include all twenty or twenty-seven measured methods, even where compact manuscript tables omit controls. The query diagnostic CSV identifies reference policies and candidate pools; the seed CSV retains request status and phase counts. Numerical ties consistently use a $10^{-12}$ tolerance. The native resource CSV distinguishes validated, memory-limited, unattempted, and synthetic cases. Thus a missing measurement, a retained incumbent, and a successful quantum simulation remain distinct events.
